\documentclass{article}
\usepackage{amsmath,amssymb,amsthm,mathtools}
\usepackage{proofread}
\newtheorem{lemma}{Lemma}
\newtheorem{theorem}[lemma]{Theorem}
\newtheorem{corollary}[lemma]{Corollary}

\title{A number and its reciprocal}
\author{Proofread demo}
\date{}

\begin{document}
\maketitle

\begin{abstract}
We show that a positive number and its reciprocal never sum to less than two, and that the sum is smallest at one.
The proof rests on an identity for the distance between the sum and two.
\end{abstract}

\section{Introduction}

This bound is the inequality between the arithmetic and the geometric mean of a positive number and its reciprocal, a classical inequality treated at length in \cite{hardy1952inequalities}.
We derive it directly from an identity, which also shows where the sum attains its minimum.
Every claim of this note is verified in Lean with Proofread.

\section{The bound}

Define $f:(0,\infty)\to\mathbb{R}$ by $f(x)\coloneqq x+\frac1x$.

\begin{lemma}\label{lem:identity}
For every $x\in(0,\infty)$, $f(x)-2=\frac{(x-1)^2}{x}$.
\end{lemma}

\begin{proof}
Let $x\in(0,\infty)$.
Then $f(x)-2=x+\frac1x-2=\frac{x^2-2x+1}{x}=\frac{(x-1)^2}{x}$.
\end{proof}

\begin{theorem}\label{thm:bound}
For every $x\in(0,\infty)$, $f(x)>2$.
\end{theorem}

\begin{proof}
Let $x\in(0,\infty)$.
Then $\frac{(x-1)^2}{x}\ge0$.
By Lemma~\ref{lem:identity}, $f(x)-2=\frac{(x-1)^2}{x}\ge0$.
\end{proof}

\begin{corollary}\label{cor:min}
For every $x\in(0,\infty)$, $f(1)\le f(x)$.
\end{corollary}

\begin{proof}
By Theorem~\ref{thm:bound}, $f(1)=2\le f(x)$ for every $x\in(0,\infty)$.
\end{proof}

\begin{thebibliography}{1}
\bibitem{hardy1952inequalities} G.~H. Hardy, J.~E. Littlewood and G.~P\'olya, \emph{Inequalities}, 2nd ed., Cambridge University Press, 1952.
\end{thebibliography}

\end{document}
